\documentclass{article}
\usepackage[backend=biber, sorting=none]{biblatex}
\usepackage{graphicx}
\addbibresource{uni.bib}

\author{Angelo Battaglia}
\title{Internship Argo Software, report 2}
\begin{document}

\maketitle

\begin{center}
\includegraphics[width=3in]{./images/polito.jpg}
\includegraphics[width=3in]{./images/og_logoargosoft.png}
\end{center}

\section{The Time Table Problem in the Italian School System}

The Time Table Problem belongs to the class of
By citing \textcite{colorni1992genetic} let's start formalizing 
the model of the Italian School System.

\subsection{Teachers}

Teachers have an eighteen hours workweek. These eighteen hours are
considered to be stored into a number of didactic units, one of which
might be considered to be composed of one or more hours. In the simulation
these units are being called U.D. (Unità Didattiche). The variable totUD is
the total of these didactic units, while minUD and maxUD represent respectively
the minimum and the maximum, of that teacher for that specific class. Therefore, 
teachers may teach one or more subjects, in two or more classes. 

\subsection{Classes}

\section{Genetic Algorithm Introduction}

\printbibliography[title = {Bibliography}]
\end{document}
